% 第4章 計量摂動とゲージ
\section{計量摂動とゲージ}
\label{ch:ch04_gauge}

\paragraph{この章で導くこと（入力→出力）}
摂動の SVT 分解を最小限示し（入力）、ゲージ自由度を理解した上で conformal
Newtonian ゲージを\emph{理由つきで}固定する（出力 S4.1）。速度・密度ゆらぎの符号
規約をここで一意に定め、第\ref{ch:ch05_einstein}--\ref{ch:ch06_boltzmann}章と
\texttt{cmbcore} の実装を整合させる。

\subsection{なぜ摂動論か}
CMB のゆらぎは $\Delta T/T\sim10^{-5}$ であり、対応する計量・密度ゆらぎも
$\lesssim10^{-4}$。よって背景＋1次摂動の線形理論が再結合期まで（小スケールの
非線形成長を除き）正当化される。本書は1次摂動のみを扱う。

\subsection{計量摂動の SVT 分解}
一般の計量摂動 $\delta g_{\mu\nu}$ は、空間回転の既約表現でスカラー(S)・
ベクトル(V)・テンソル(T)に分解できる。1次では三者は独立に発展し、ベクトルは
減衰、テンソルは原始重力波（本書の主筋外、\textbf{発展}）。スカラーは4つの自由度を
持つが、ゲージ固定で物理的な2つに減る。

\begin{advanced}
ベクトル摂動は能動的源がなければ $\propto a^{-2}$ で減衰する。テンソル摂動（重力波）は
B モード偏光を生むが、TT スペクトルへの寄与は無視できる。本書はスカラーのみ扱う。
\end{advanced}

\subsection{ゲージ変換（導出契約 S4.1）}
無限小座標変換 $x^\mu\to x^\mu+\xi^\mu$ のもとで計量・物質の摂動は
\begin{equation}
  \delta g_{\mu\nu}\to\delta g_{\mu\nu}-\nabla_\mu\xi_\nu-\nabla_\nu\xi_\mu,\qquad
  \delta\rho\to\delta\rho-\dot{\bar\rho}\,\xi^0,
\end{equation}
と変換する（$\xi^\mu=(\xi^0,\xi^i)$、$\xi^i=\partial^i\xi$ がスカラー部）。
すなわち\textbf{密度ゆらぎはゲージ依存}である。

\begin{intuition}
\textbf{ゲージ依存の具体例}。時間スライス（同時刻面）を $\xi^0$ だけずらすと、各点の
背景密度が $\dot{\bar\rho}\,\xi^0$ だけ変わるので、観測される $\delta\rho$ もその分動く。
極端には「密度が一様な面」を選べば $\delta\rho=0$（$\xi^0=\delta\rho/\dot{\bar\rho}$）に
できてしまう。超地平線では $\delta\rho/\bar\rho$ がゲージで $O(1)$ も変わりうるため、
「どの面で測った密度か」を宣言せずに摂動を語るのは無意味である。一方、計量の曲率に
基づく $\mathcal R$（下記）はゲージ不変で、物理的な「ゆらぎの大きさ」を担う。
\end{intuition}

物理は変換で不変なゲージ不変量、または残留自由度のないゲージで固定して扱う。本書は
後者（conformal Newtonian）を採る。

\subsection{conformal Newtonian ゲージの固定}
スカラー摂動の2自由度を使い、計量を
\begin{equation}
  ds^2=-(1+2\Psi)\,c^2dt^2+a^2(t)(1+2\Phi)\,\delta_{ij}\,dx^idx^j,
  \label{eq:metric}
\end{equation}
に固定する（Dodelson/Callin 規約、`00` §3.1）。この形は非対角・非等方成分を持たず
残留ゲージ自由度がない。$\Psi$ はニュートンポテンシャル（時間成分）、$\Phi$ は空間
曲率摂動で、異方性応力が無視できるとき $\Phi\simeq-\Psi$（第\ref{ch:ch05_einstein}章）。

\begin{numreality}
Ma \& Bertschinger の $(\psi,\phi)$ とは\textbf{符号規約が異なる}（付録A の対応表）。
\texttt{cmbcore} は式\eqref{eq:metric} の $(\Psi,\Phi)$ をそのまま用いる。
\end{numreality}

\subsection{ゲージ不変量 $\mathcal R$}
共動曲率ゆらぎ $\mathcal R$ は超地平線で保存する量で（第\ref{ch:ch07_initial}章）、
原始スペクトル $\mathcal P_\mathcal R$ の規格化に使う。Newtonian ゲージ変数では
$\mathcal R=\Phi-\frac{\mathcal H}{k}v$ 型（本書符号）で表され、断熱初期条件を
$\mathcal R=1$ に規格化する（第\ref{ch:ch11_cl}章の規格化監査）。

\subsection{流体変数の定義}
密度コントラスト $\delta_i\equiv\delta\rho_i/\bar\rho_i$、速度 $v_i$ を
$T^0{}_i=(\bar\rho+\bar p)v_i$、スカラー速度 $v$ を $v_i=-\frac{ik_i}{k}v$ で定義する
（`03_numerics_spec.md` と一致）。光子は輝度 $\Theta$ の多重極で表し、
$v_\gamma=-3\Theta_1$ の関係を持つ（第\ref{ch:ch06_boltzmann}章）。

\paragraph{まとめ}
SVT 分解からスカラー2自由度を取り出し、conformal Newtonian ゲージ\eqref{eq:metric} を
固定、$\delta_i,v_i,\mathcal R$ の符号規約を確定した。

\paragraph{演習}
(1) $\delta\rho$ のゲージ変換則から一様密度ゲージで $\delta\rho=0$ を示せ。
(2) 式\eqref{eq:metric} に残留ゲージ自由度がないことを確認せよ。
(3) Ma\&Bertschinger 規約との符号対応を付録A で確かめよ。
